\documentclass{article}
\usepackage[margin=2cm]{geometry}
\usepackage{amsmath}
\usepackage{graphicx}
\usepackage{booktabs}
\usepackage[utf8]{inputenc}
\usepackage{ragged2e}
\setlength{\emergencystretch}{3em}
\begin{document}
\sloppy

The expectation values of the operator Og3 and the trace of the energy-momentum tensor
are independent of . As a check, we note that the four one-point functions satisfy the
following operator relation, which is associated to the violation of conformal invariance by.
non-zero classical beta functions,


\begin{equation}
\langle T ^ { i } { } _ { i } \rangle = - \sum _ { \mathcal { O } } ( d - \Delta _ { \mathcal { O } } ) ~ \phi _ { \mathcal { O } } \left < \mathcal { O } \right > \qquad \qquad ( 1 )
\end{equation}


\subsection*{0.0.1 Derivative of the free energy}


Following , we may now compute the derivative of F with respect to a as follows. First we
note that


\begin{equation}
\frac { d F } { d \alpha } = \frac { d S _ { \mathrm { s e n } } } { d \alpha } = \operatorname * { l i m } _ { \epsilon \rightarrow 0 } \int d ^ { 5 } x \sum _ { \mathrm { h e l d } \cdot \Phi } \frac { \delta \left ( \sqrt { \gamma } \mathcal { L } _ { \mathrm { r e n } } \right ) } { \delta \Phi } ~ \frac { d \Phi } { d \alpha } \Bigg | _ { z = : } \qquad ( 2 )
\end{equation}


In our case, the terms appearing in the sum over fields are.


\begin{equation}
{ \begin{array} { c c c } { { \frac { \delta \left ( \sqrt { \gamma } \mathcal { L } _ { \mathrm { r e n } } \right ) ^ { \! \! \! \! \! / } } { \delta \sigma } } } & { { \stackrel { \! \! \! \! \! / } { \delta } } } \\ { { \frac { \delta \left ( \sqrt { \gamma } \mathcal { L } _ { \mathrm { r e n } } \right ) } { \delta \sigma ^ { 3 } } } + \dots ~ }}\end{array}
\end{equation}


The dots represent terms of strictly lower order in e. Furthermore, from the form of the UV
asymptotic expansions , we have


\begin{equation}
\begin{array} { c c } { } & { } \\ { } & { \frac { d \sigma } { d \alpha } = \underbrace { 3 } _ { \mathrm { t o } } \alpha \epsilon ^ { 2 } + O ( \epsilon ^ { 3 } ) \qquad \qquad \qquad ~ }\end{array}
\end{equation}


Combining the pieces , with the results for the one-point functions in , we find that the.
contribution of the metric in is suppressed by e? compared to other terms. The derivative of
the free energy is then.


\begin{equation}
{ \begin{array} { c c } { { \frac { d F } { d \alpha } } = \operatorname * { l i m } _ { \epsilon \rightarrow 0 } \int d ^ { 5 } x \sqrt { \gamma } ~ \epsilon ^ { 6 } \left [ { \frac { 3 } { 2 } } \beta e ^ { - f _ { k } } + { \frac { 1 } { 2 } } \beta e ^ { - f _ { k } } \left ( 1 - \alpha { \frac { d f _ { k } } { d \alpha } } \right ) + O ( \epsilon ) \right ] } \\ { { } } & { { } } \\ { { = \mathrm { v o l } _ { 0 } \left ( S}}}\end{array}
\end{equation}


where volo(S) = 3 is the volume of a unit S5. The  dependence in the one-point functions
cancels out, consistent with the fact that F itself is independent of . We thus obtain the
final result


\begin{equation}
\frac { d F } { d \alpha } = \frac { \pi ^ { 2 } } { 8 ~ G _ { 6 } } \beta ~ e ^ { 4 f _ { k } } \left ( 4 - \alpha \frac { d f _ { k } } { d \alpha } \right ) \qquad ( \mathbf { 6 } )
\end{equation}

\end{document}
